\section{问题重述与总体分析}

波浪能转换装置由浮子、振子、弹簧和阻尼器构成。波浪激励力使浮子产生垂荡或纵摇运动，浮子与内部振子之间的相对运动驱动 PTO 阻尼器做功，从而实现能量输出。题目要求在不同运动自由度和不同阻尼形式下，建立动力学模型并求解运动响应，进一步优化阻尼参数以最大化平均输出功率。

本文将四个子问题分为两类。问题一、二仅考虑垂荡运动，系统可抽象为浮子和振子的二自由度受迫耦合振动模型。问题一在给定阻尼参数下求运动响应；问题二在相同动力学框架下改变 PTO 阻尼参数，以稳态平均功率作为优化目标。问题三、四进一步考虑纵摇，系统自由度增至四个，几何关系导致方程中出现三角函数、变转动惯量和惯性耦合项，因此采用数值积分求解。问题四在问题三模型基础上同时优化垂荡线性阻尼和旋转阻尼。

本文采用的总体技术路线为：首先依据附件参数建立受力与力矩平衡方程；其次将高阶微分方程转化为一阶状态方程并用数值积分求解；最后通过稳态区间功率积分评价输出功率，并结合粗扫描与局部优化寻找最优阻尼参数。

\section{数据来源与参数处理}

结构参数来自附件4，包括浮子质量 $m_b=4866\,\mathrm{kg}$、振子质量 $m_z=2433\,\mathrm{kg}$、弹簧刚度 $k=80000\,\mathrm{N/m}$、海水密度 $\rho=1025\,\mathrm{kg/m^3}$、重力加速度 $g=9.8\,\mathrm{m/s^2}$ 等。水动力参数来自附件3，不同问题对应不同波浪圆频率、附加质量、兴波阻尼和激励力幅值。

浮子垂荡静水恢复力系数由水线面积给出：
\begin{equation}
  K_h=\rho g\pi r_b^2 .
  \label{eq:kh}
\end{equation}
在考虑纵摇时，静水恢复力矩系数取附件给定值 $K_\theta=8890.7\,\mathrm{N\cdot m}$。浮子纵摇转动惯量按照圆柱壳与圆锥壳表面积分配质量的近似公式计算；振子绕横向中心轴的转动惯量为
\begin{equation}
  I_z=\frac{1}{12}m_z(3r_z^2+h_z^2).
  \label{eq:iz}
\end{equation}
